% !TeX root = ../papers/nc-p3.tex
\section{楕円座標端の全通常点での正則性}
\label{sec:rank-fifteen-coordinate-regularity}
\begin{theorem}\label{thm:rank-fifteen-coordinate-regularity}
標数零で $n=\pm1$ の両表示の全通常な階数15窓は標準Koszul左右AS正則である。
各始点の道空間は四本の四次中心元の生成する多項式環上で階数256の自由加群になる。
\end{theorem}
\begin{proof}
二位置移動と倍率反転により第一表示の $n=1$ を取る。
定理\ref{thm:rank-fifteen-coordinate-ends}の楕円曲線上では四次までの経路、
長さ5の経路、四次中心は既に一定次元である。
その関数体 $\mathbb Q(t,Y)$ 上で、中心の文字0の三自由座標の第一と第三、
文字1の二座標の両方を選ぶ。これら四本で割った商について、定理
\ref{thm:rank-fourteen-all-ordinary}の右への矢の追加と関係挿入を行う。
四本は全位置に共通の中心切断である。

\path{computations/verify_rank_fifteen_coordinate_regularity.py} は逆元を取る順を
疎な行と次数の低い係数を優先するように選ぶ。各係数を $(A+BY)/D$ の形で既約化し、
逆元を取る全ての分子のノルム $A^2-(-t^3+14t^2-t)B^2$ の既約因子を記録する。
初期座標での除算も含む。始点0,1の各々で長さ0から13までの商次元の上界は
\[
 1,4,10,20,31,40,44,40,31,20,10,4,1,0
\]
となる。各段で消去に用いた行の関係式だけでも、残した語が商を張る。
二位置移動の同型で他の二始点にも良い四本が存在する。四つの始点で共通の四本は、
中心の同じ係数空間で四つの非空開集合の交叉を取って得られる。

全除算に現れる相異なる既約因子は464個である。前定理の同時反転は楕円上で
$(t,Y)\mapsto(1/t,-Y/t^2)$ となる。ノルムは両方の $Y$ を覆うので、
$t^{\deg g}g(1/t)$ も記録にある因子だけを残せばよい。
$t=0$ を除くこの共通部分は69因子で、非通常因子を除き相反な因子を組にすると55軌道になる。
コードは各代表の単根と $Y$ の二つの相異なる剰余値を指定し、その両方で
経路・中心の次元と全四始点の長さ13までの商を検証する。
前節の座標端の場合と同じく、単根での成功と共役により各軌道全体を覆う。
一般点の道次元の下界とここでの中心商の上界から自由加群の上下界が一致する。
定理\ref{thm:rank-fourteen-all-ordinary}の自由性と左右の解消の議論で結論が従う。
\end{proof}

\begin{proposition}\label{prop:endpoint-moduli-normalizations}
座標端から得られる一般添字保持同型類を与える三つの像曲線は相異なる。
その滑らかな射影正規化は $j=35152/9$ の楕円曲線と二本の射影直線である。
\end{proposition}
\begin{proof}
定理\ref{thm:endpoint-generic-fibers}を与える不変量写像の各成分上の一般次数は2であり、
その関数体は同時反転による不変部分体に一致する。楕円上では
$s=t+1/t$、$v=(Y/t)(t-1/t)$ が同時反転で不変である。
\[
 v^2=(14-s)(s^2-4).
\]
$t^2-st+1=0$ と $v$ から元の関数体を二次で復元できるので、これが求める商である。
三根 $-2,2,14$ のLegendreパラメータは4で、記載の $j$ を得る。
有理成分では $s=t-1/t$ が不変部分体を生成する。

楕円と有理の像の相違は前定理で得た。有理曲線の二つの符号成分の間では、
$i$ 成分の $t=2$ の $J_2,J_3$ を $-i$ 成分で取る二式の最大公約式は
$(t+2)(t+3)(t-1/2)(t-1/3)$ で、無限遠にも解はない。
\path{computations/verify_rank_fifteen_endpoint_orbits.py} は
$109$ 上で $i=33$ を選び、元の点で $L=46$、相手の四候補全てで $L=63$ を得る。
各点は通常な階数15で全分母は単元である。従って元の不変量の値は相手の完備曲線の
像に属さず、両像も異なる。
\end{proof}
